\section{Termination and the log canonical reduction}\label{sec:completion}

We first show that a surface in the positive exceptional locus of a
downstairs diagram is detected by the integer difficulty.  The calculation
uses the weight denominator \(N\) chosen for the particular strong gklt
sequence under consideration.  We then prove
termination for strong gklt models, pass to elementary gdlt programs, and
lift the arbitrary diagrams of Theorem~\ref{thm:main} to those programs.

\begin{lemma}[A positive exceptional surface]\label{lem:positive-surface}
Use the truncated terminal chain of
Propositions~\ref{prop:terminal-junctions} and
\ref{prop:surface-finiteness}, with the integer \(N\), weights, and constant
defaults chosen in Section~\ref{sec:difficulty}.  If an irreducible compact
analytic surface \(V\) is contained in \(L_i^+\subset X_{i+1}\), there is
a global place \(E\), exceptional over both junctions and every model of
their connector, such that
\[
 a_{i+1}(E)\in(1,2]\cap N^{-1}\Z,
 \qquad w_i(E)\geq w_{i+1}(E)+1.
\]
Here \(w_i\) and \(w_{i+1}\) are the weights at the two crepant junctions.
Consequently
\[
 \mathfrak D(Y_i,\Theta_i+\M)
 >\mathfrak D(Y_{i+1},\Theta_{i+1}+\M).
\]
\end{lemma}

\begin{proof}
By Proposition~\ref{prop:surface-finiteness}, the model \(X_{i+1}\) is
smooth near the general point of \(V\), the junction morphism is an
isomorphism there, and the generic blowup of \(V\) gives a unique global
place \(E\) with
\[
 a^+:=a_{i+1}(E)>1.
\]
This place is exceptional over \(X_{i+1}\) and does not belong to the
stabilized extracted set \(\mathcal I\).  The blowup of the coherent ideal
of \(V\) is projective, so Lemma~\ref{lem:common-projective-models} gives a
smooth common model \(W\) carrying both \(E\) and the fixed nef datum,
with a projective morphism \(r:W\to X_{i+1}\).  By Lemma~\ref{lem:nef-defect},
\[
 J=r^*M_{i+1}-M_W\geq0.
\]
At the generic smooth codimension-two center, the ordinary log discrepancy
of the blowup is two.  The discrepancy formula is therefore
\begin{equation}\label{eq:positive-detector-discrepancy}
 a^+=2-\ord_E(B_{i+1})-\operatorname{coeff}_E(J).
\end{equation}
Both subtracted terms are nonnegative.  Thus \(a^+\leq2\).

We verify the denominator in this calculation.  By the choice in
Section~\ref{sec:difficulty}, \(N\) clears the coefficients of the initial
boundary and the Cartier denominator of the higher nef divisor of this
strong gklt sequence.  Since the downstairs boundaries are strict
transforms, \(NB_{i+1}\) is an integral Weil divisor.  Also \(NM_{i+1}\)
is an integral Weil divisor: it is the pushforward of the integral Cartier
divisor \(NM_W\) from a common carrier.  On the smooth open neighborhood of
the general point of \(V\), these two integral Weil divisors are Cartier.
Their orders at \(E\) are consequently integral after multiplication by
\(N\).  In particular,
\[
 N\ord_E(B_{i+1})\in\Z,
 \qquad
 N\operatorname{coeff}_E(J)
 =\ord_E(Nr^*M_{i+1})-\operatorname{coeff}_E(NM_W)\in\Z.
\]
Only Cartierness on this smooth open has been used.  Equation
\eqref{eq:positive-detector-discrepancy} gives
\(a^+\in(1,2]\cap N^{-1}\Z\).

Put \(a^-=a_i(E)\).  The strict part of
Lemma~\ref{lem:small-comparison} applies because
\(\cent_{X_{i+1}}(E)=V\subset L_i^+\), and gives \(a^-<a^+\).  The same
lemma says that \(E\) is exceptional over both downstairs models.  Since
it does not belong to \(\mathcal I\), it is not an extracted prime on either
junction; it is therefore exceptional over both.  All the intervening
steps are small, so it stays exceptional on every model of the connector.

Crepancy identifies the discrepancies at the junctions with \(a^-\) and
\(a^+\).  Write \(k=N(2-a^+)\), a nonnegative integer.  Since
\(a^-<a^+\leq2\), the ceiling weights satisfy
\[
 w_i(E)=\left\lceil N(2-a^-)\right\rceil
 \geq k+1,
 \qquad
 w_{i+1}(E)=\left\lceil N(2-a^+)\right\rceil=k.
\]
This includes \(a^+=2\), when \(k=0\).  No denominator assertion for
\(a^-\) is needed.  The connector consists of one coefficient change and
a finite, possibly empty, string of small steps.  Its weights are
nonincreasing, so the strict decrease between its ends occurs in at least
one move.  The place \(E\) is exceptional throughout that move.
Proposition~\ref{prop:difficulty-change} makes the difficulty decrease
strictly there and nonincrease in every other move, proving the claim.
\end{proof}

\begin{theorem}[Termination on strong gklt models]\label{thm:strong-gklt}
Every sequence satisfying the hypotheses of Theorem~\ref{thm:main} for
which all the models \(X_i\) are globally strongly \(\Q\)-factorial and
the initial pair \((X_0,B_0+\M)\) is gklt is finite.  The small diagrams
may be chosen arbitrarily subject to the hypotheses of that theorem.
\end{theorem}

\begin{proof}
Suppose such a sequence were infinite.  Apply
Proposition~\ref{prop:terminal-junctions}, then
Proposition~\ref{prop:surface-finiteness}, and use the further truncation
and denominator choice of Section~\ref{sec:difficulty}.  At every junction
the difficulty is a nonnegative integer by
Proposition~\ref{prop:difficulty-lower}.  It is nonincreasing from one
junction to the next by Proposition~\ref{prop:difficulty-change}.  By
Lemma~\ref{lem:positive-surface}, it drops by at least one whenever the
positive exceptional locus contains a surface.  There can be only finitely
many such indices.  This argument uses the lower bound only at junctions;
it requires no lower bound on the intermediate models.

Pass to the remaining tail.  There \(\dim L_i^+<2\).  The dimension
inequality in Lemma~\ref{lem:small-comparison} forces \(\dim L_i=2\), so
\(L_i\) contains an irreducible compact analytic surface.  By
Lemma~\ref{lem:cycle-loss},
\[
 c_2(X_i)>c_2(X_{i+1})
\]
at every step of this tail.  These are nonnegative integers, a
contradiction.
\end{proof}

\begin{theorem}[Termination of elementary gdlt programs]\label{thm:gdlt}
Let
\[
 (T_0,B_0+\M)\dashrightarrow(T_1,B_1+\M)
 \dashrightarrow(T_2,B_2+\M)\dashrightarrow\cdots
\]
be a sequence of elementary steps for rational gdlt pairs on strong normal irreducible
compact K\"ahler fourfolds.  Assume that the boundaries are effective
rational divisors of finite support transported by the steps, that the
fixed b-divisor is represented by an analytically nef \(\Q\)-Cartier divisor
on a projective bimeromorphic normal compact K\"ahler model of \(T_0\), and
that the actual adjoints are \(\Q\)-Cartier with compatible canonical
representatives.  Then the sequence is finite.  The contraction bases and
the choices of elementary steps may vary.
\end{theorem}

\begin{proof}
An elementary program extracts no prime divisors.  Lemma~\ref{lem:cycle-loss} with \(k=3\)
therefore permits only finitely many divisorial steps.  If the sequence were
infinite, we could discard a finite prefix containing them and work on a
tail of small steps.
The rational b-line realization described in
Section~\ref{sec:preliminaries} has the same discrepancies and fixed higher
nef data; common projective carriers for finite portions carry their
analytically nef pullbacks.  Thus Theorem~\ref{thm:special-termination}
applies.  After a further tail, both exceptional loci of every step avoid
every generalized log canonical center on their respective sides.

Set \(H_i=\lfloor B_i\rfloor\).  Since the pair is glc and \(B_i\) is
effective, its coefficients belong to \([0,1]\), so \(H_i\) is the reduced
sum of its coefficient-one components.  Each such component is itself a
generalized log canonical center.  Hence \(\Supp H_i\) is disjoint from
the exceptional loci adjacent to \(T_i\) on this tail.  The divisors
\(H_i\) are effective and \(\Q\)-Cartier on the strong models, and they
are strict transforms of each other across the small steps.

Remove the whole floor and write
\[
 \widehat B_i=B_i-H_i,
 \qquad \widehat D_i=K_{T_i}+\widehat B_i+M_{T_i}=D_i-H_i.
\]
The new boundary is effective.  For every global place \(E\), comparison
of the two crepant equations gives
\begin{equation}\label{eq:full-floor-deletion}
 a(E;T_i,\widehat B_i+\M)
 =a(E;T_i,B_i+\M)+\ord_E(H_i).
\end{equation}
The added order is nonnegative.  If the old discrepancy is zero, the gdlt
definition says that the general point of its center is a stratum of the
coefficient-one boundary.  Taking its closure places the entire center in
\(\Supp H_i\).  The order of an effective \(\Q\)-Cartier divisor at a
place whose center lies in its support is strictly positive.  Thus every
former zero discrepancy becomes positive, and every former positive
discrepancy stays positive.  The pairs \((T_i,\widehat B_i+\M)\) are gklt.

We check the two ample signs after this full deletion.  For a step write
\[
 T_i\xrightarrow{f}Z\xleftarrow{f^+}T_{i+1},
 \qquad L=\Exc(f),\quad L^+=\Exc(f^+).
\]
Lemma~\ref{lem:small-comparison} gives a common exceptional image \(A\)
with \(L=f^{-1}(A)\), \(L^+=(f^+)^{-1}(A)\), and both maps isomorphisms
off \(A\).  Properness makes \(f(\Supp H_i)\) closed.  It is disjoint from
\(A\), because \(\Supp H_i\cap L=\varnothing\).  On the base open
\(Z\setminus f(\Supp H_i)\), which contains \(A\), the divisor \(H_i\)
vanishes, so \(-\widehat D_i=-D_i+H_i\) has the original relative ample
sign.  On \(Z\setminus A\) the map \(f\) is an isomorphism, and every line
bundle is relatively ample for an isomorphism.  These two opens cover
\(Z\).  After clearing a common Cartier multiple, relative ampleness is
local on the base, so \(-\widehat D_i\) is \(f\)-ample.  The identical
argument with \(f^+\) and \(H_{i+1}\) proves that
\(\widehat D_{i+1}=D_{i+1}-H_{i+1}\) is \(f^+\)-ample.

Reindex this floor-deleted tail at its first model and regard it as a
sequence of the arbitrary small diagrams in Theorem~\ref{thm:strong-gklt}.
The model and diagram hypotheses are inherited; the new boundaries are
effective and related by strict transform, the actual adjoints are
\(\Q\)-Cartier, and the two ample signs have just been checked.  A common
carrier for the finite preceding portion, projective over the new first
model, supplies the fixed analytically nef divisor.  In this application
the detector integer is chosen afresh for the new initial boundary and
nef carrier.  The theorem makes the tail finite, contradicting its
construction.
\end{proof}

To lift a prescribed downstairs diagram, we run gdlt continuation over
its base.  Theorem~\ref{thm:gdlt} forces the chosen run to reach a nef
adjoint, and Lemma~\ref{lem:nef-ample} gives a projective crepant morphism
from that endpoint to the prescribed positive model.  A negative
multisection ensures that the run contains a step.

\begin{lemma}[A nonempty gdlt lift]\label{lem:nonempty-lift}
Let
\[
 X\xrightarrow{\ f\ }Z\xleftarrow{\ f^+\ }X^+
\]
be any one of the small diagrams allowed in Theorem~\ref{thm:main}, with
normal globally Weil \(\Q\)-factorial compact K\"ahler fourfolds on the two
sides, effective rational boundaries \(B,B^+\) related by strict transform,
the same rational nef b-divisor, and the compatible actual adjoints
\(D_X,D_{X^+}\).  Assume \((X,B+\M)\) is glc, and represent the b-divisor
by an analytically nef \(\Q\)-Cartier divisor on a projective carrier over
\(X\).  Suppose
\[
 h:(Y,\Theta+\M)\longrightarrow(X,B+\M)
\]
is a projective crepant gdlt modification as in
Proposition~\ref{prop:gdlt-modification}; in particular \(Y\) is strong,
\(\Theta\geq0\), and every \(h\)-exceptional prime has coefficient one.
Then there is a finite sequence of elementary negative gdlt steps over
\(Z\), of length \(m\geq1\),
\[
 (Y,\Theta+\M)=(Y_0,\Theta_0+\M)
 \dashrightarrow\cdots\dashrightarrow(Y_m,\Theta_m+\M),
\]
and a projective bimeromorphic morphism \(h^+:Y_m\to X^+\) such that
\[
 K_{Y_m}+\Theta_m+M_{Y_m}=(h^+)^*D_{X^+}.
\]
The end is a strong normal compact K\"ahler gdlt model with effective
boundary, and every \(h^+\)-exceptional prime has coefficient one in
\(\Theta_m\).  Thus it is again a crepant gdlt modification of the same
type.
\end{lemma}

\begin{proof}
The composite \(Y\to Z\) is projective bimeromorphic to a normal compact
K\"ahler base.  On a projective common carrier, the fixed analytically nef
divisor is nef over \(Z\).  Thus the continuation assertion of
Proposition~\ref{prop:relative-program} applies to this gdlt pair and to
every non-nef finite prefix produced from it.

The initial adjoint \(D_Y=h^*D_X\) is not nef over \(Z\).  Indeed, the
nonisomorphic projective contraction \(f\) has a positive-dimensional
fiber, and projective hyperplane cuts give a compact irreducible curve
\(C\) in that fiber.  The negative ample sign gives \(D_X\cdot C<0\).
There is a compact irreducible curve \(C'\subset Y\) mapping onto \(C\).
For completeness, normalize \(C\), take an irreducible component of its
base change under \(h\) which dominates that normalization, and write
\(r\) for its generic fiber dimension.  The normalization of \(C\) is a
projective compact curve, so this component, being projective over it, is
projective.  A sufficiently high very ample system on the component has
\(r\) general members whose intersection with a generic fiber is a
nonempty zero-dimensional cycle.  Their intersection therefore has a
compact curve component dominating the base.  Its image in \(Y\) is such
a \(C'\).  The induced map of normalized curves has a positive finite
degree.  The projection formula now gives
\[
 D_Y\cdot C'=\deg(C'/C)\,D_X\cdot C<0.
\]
In particular, any run reaching nefness must contain at least one step.

Choose an elementary negative step whenever the current adjoint is not
nef over \(Z\).  The relative continuation assertion keeps all outputs
projective over \(Z\), strong normal compact K\"ahler, and gdlt with the
transported data.  All models in the run are irreducible: \(Y\) is
bimeromorphic to the irreducible \(X\), and every elementary step is
bimeromorphic.  If this process never reached a nef adjoint, it would
give an infinite elementary gdlt sequence with the fixed analytically nef
b-divisor.  This contradicts Theorem~\ref{thm:gdlt}.  We obtain a finite
run of length \(m\geq1\) with end adjoint \(D_{Y_m}\) nef over \(Z\).

Take a smooth common model \(W\) with projective maps to the models in the
run, to \(X^+\), and to the nef carrier, all over \(Z\).  Write \(P_0,P,P_+\) for the pullbacks
of \(D_Y,D_{Y_m},D_{X^+}\), respectively.  By
Lemma~\ref{lem:program-comparison}, the finite run gives
\[
 P_0=P+G,\qquad G\geq0\text{ exceptional over }Y_m.
\]
Initial crepancy identifies \(P_0\) with the pullback of \(D_X\).
Lemma~\ref{lem:small-comparison}, pulled to \(W\), gives the second
comparison
\[
 P_0=P_++F,\qquad F\geq0\text{ exceptional over }X^+.
\]
The first end is nef over \(Z\), and \(D_{X^+}\) is ample over \(Z\).
Lemma~\ref{lem:nef-ample} therefore gives an actual equality \(P=P_+\)
and a projective morphism \(h^+:Y_m\to X^+\) over \(Z\) with
\(D_{Y_m}=(h^+)^*D_{X^+}\).

The end is strong, gdlt, and has effective boundary by the relative
program.  It remains to check its exceptional primes.  Every prime on
\(Y_m\) has a strict-transform prime on \(Y\), because an elementary
program extracts no prime divisors.  Let \(R\) be exceptional for
\(h^+\), and let \(R_0\) be its prime on \(Y\).  If \(R_0\) were not
exceptional for \(h\), it would represent a prime on \(X\).  Smallness
downstairs identifies that prime with a prime on \(X^+\).  The place of
\(R\) would then have a divisorial center on \(X^+\), contradicting its
exceptionality.  Hence \(R_0\) is \(h\)-exceptional and has coefficient
one in \(\Theta\).  Boundary coefficients of surviving primes are
unchanged by an elementary program, so \(R\) has coefficient one in
\(\Theta_m\).  The displayed crepant equality also identifies its
discrepancy over \(X^+\) as zero.  This proves the lemma.
\end{proof}

\begin{proof}[Proof of Theorem~\ref{thm:main}]
Assume that the sequence in the theorem is infinite.  By
Lemma~\ref{lem:small-comparison}, every global discrepancy is nondecreasing
along it.  Hence every \((X_i,B_i+\M)\) is glc.  The rational b-line
realization of the initial data satisfies
Proposition~\ref{prop:gdlt-modification}: \(X_0\) is a normal compact
K\"ahler fourfold, the boundary is effective and rational, the nef datum
has a projective carrier, and the adjoint is rationally invertible.  The
proposition gives a projective crepant gdlt modification
\[
 h_0:(Y_0,\Theta_0+\M)\longrightarrow(X_0,B_0+\M)
\]
with a strong normal compact K\"ahler source and only coefficient-one
exceptional primes.  It requires no strong factoriality of \(X_0\).

Apply Lemma~\ref{lem:nonempty-lift} to the first downstairs diagram and
\(h_0\).  Its endpoint is a crepant gdlt modification of the same type over
\(X_1\), so it is the input for the next application.  At every finite
stage a common projective carrier supplies the fixed analytically nef
data over the current downstairs model.  Induction therefore lifts every
downstairs diagram.  Each lifted string is finite and contains at least one
elementary step, and its endpoint is used as the literal starting model of
the next string.  Their concatenation is an infinite sequence of elementary
gdlt steps on strong normal irreducible compact K\"ahler fourfolds, with the same
b-divisor throughout.  Theorem~\ref{thm:gdlt} allows the bases to vary and
rules out this sequence.  The contradiction proves
Theorem~\ref{thm:main}.
\end{proof}
